\chapter[Introdução]{Introdução}

\section{Contextualização do Tema}
A manutenção e operação eficiente de infraestruturas críticas, como subestações de energia elétrica, são fundamentais para garantir a estabilidade e a segurança do fornecimento de energia à sociedade \cite{suitelec2023}. Essas instalações, frequentemente complexas e localizadas em áreas remotas ou de difícil acesso, demandam inspeções regulares para identificar potenciais falhas, desgastes ou condições anormais que possam comprometer seu funcionamento ou representar riscos. Tradicionalmente, essas inspeções são realizadas por equipes humanas, um processo que, embora essencial, pode ser custoso, demorado, sujeito a erros humanos e, em certos casos, expor os trabalhadores a ambientes perigosos \cite{Li2025SafetyRisk}.
Nas últimas décadas, a robótica móvel autônoma emergiu como uma solução promissora para automatizar tarefas de inspeção em diversos setores industriais \cite{Wang2024Advancements}. A aplicação de robôs em subestações elétricas oferece o potencial de aumentar a frequência e a consistência das inspeções, reduzir custos operacionais, minimizar a exposição humana a riscos e coletar dados de forma mais sistemática e detalhada \cite{RobotnikSubstationBenefits}. Equipados com uma variedade de sensores, como câmeras visuais e térmicas, \ac{LiDAR}, microfones e detectores de gás, os robôs podem monitorar equipamentos, detectar anomalias como superaquecimento, descargas parciais, vazamentos ou danos estruturais e navegar autonomamente pelo ambiente complexo da subestação.
O desenvolvimento de sistemas de navegação autônoma para esses robôs é um campo de pesquisa ativo. Abordagens baseadas em \ac{SLAM} permitem que os robôs construam mapas do ambiente e se localizem dentro deles simultaneamente \cite{Bosche2018SLAM}. Modelos tridimensionais pré-existentes, muitas vezes gerados por técnicas como fotogrametria ou escaneamento a laser – áreas de expertise do grupo de pesquisa onde este trabalho se insere –, também podem ser utilizados para auxiliar na navegação e no planejamento de rotas. Além disso, a integração com tecnologias como \ac{RV} e \ac{RA} abre novas possibilidades para a operação remota e a visualização de dados coletados pelos robôs \cite{Zhang2023ARVR}.
Contudo, apesar dos avanços significativos, muitas das soluções atuais de inspeção robótica em subestações ainda operam predominantemente com base em rotas pré-definidas ou otimizadas para cobertura de área \cite{TJH2b2023Trends}. Essa abordagem, embora útil para tarefas de monitoramento rotineiro (ronda), pode não ser a mais eficiente ou eficaz para uma inspeção verdadeiramente investigativa. Eventos anômalos ou condições que exigem atenção imediata podem ocorrer fora da rota planejada ou exigir uma resposta dinâmica que vai além do simples seguimento de um caminho. A capacidade de um robô não apenas seguir uma rota, mas perceber ativamente seu entorno, interpretar múltiplos estímulos sensoriais de forma integrada e reagir inteligentemente a eventos inesperados, direcionando sua atenção para onde ela é mais necessária, representa um desafio e uma oportunidade significativa para a próxima geração de sistemas de inspeção autônoma.


\section{Problema de Pesquisa}
Apesar do potencial da robótica autônoma, a implementação eficaz de sistemas de inspeção em subestações elétricas enfrenta desafios significativos que limitam o desempenho e a confiabilidade das abordagens atuais, majoritariamente baseadas em rotas pré-definidas \cite{TJH2b2023Trends}. A simples execução de rondas programadas, embora útil para cobertura sistemática, falha em abordar a necessidade de uma inspeção dinâmica e focada em eventos anômalos que podem surgir inesperadamente no ambiente complexo e muitas vezes hostil de uma subestação \cite{Wang2024Advancements}. A verdadeira inteligência em inspeção reside na capacidade do robô de perceber, interpretar e reagir adequadamente a estímulos ambientais em tempo real, adaptando sua trajetória e foco de inspeção conforme necessário.
A navegação autônoma robusta neste ambiente é particularmente desafiadora devido a uma confluência de fatores ambientais e operacionais. Primeiramente, a presença de altos níveis de interferência eletromagnética ou \ac{EMI}, inerente às operações de alta tensão, pode afetar severamente o funcionamento de sensores eletrônicos e sistemas de comunicação do robô, comprometendo a aquisição de dados e a própria capacidade de navegação \cite{Liu2016EMC,tdworld2024Challenges,Wang2024CombinedNav}. Sensores de localização como GPS podem ser indisponíveis ou imprecisos, e até mesmo sensores como \ac{LiDAR} e câmeras podem sofrer interferências ou degradação de desempenho.
Adicionalmente, o ambiente físico das subestações impõe obstáculos consideráveis. A abundância de estruturas metálicas (equipamentos, grades, postes) pode causar múltiplas reflexões e oclusões para sensores baseados em ondas, como \ac{LiDAR} e radar, além de dificultar a interpretação de dados visuais \cite{Wang2024CombinedNav, Hou2023ChallengesSolutions}. A variabilidade extrema nas condições de iluminação, alternando entre luz solar direta, sombras profundas e iluminação artificial noturna, representa um desafio significativo para algoritmos de visão computacional. A necessidade de operar tanto em ambientes internos quanto externos adiciona outra camada de complexidade, exigindo que o sistema de navegação se adapte a diferentes tipos de terreno, condições atmosféricas e características espaciais \cite{Hussain2020ReviewChallenges}. A presença de obstáculos fixos e potencialmente dinâmicos (veículos de manutenção, pessoal) também requer capacidades avançadas de detecção e desvio.
Portanto, o problema central de pesquisa abordado nesta tese reside na \textbf{limitação dos atuais \textit{frameworks} de navegação robótica para inspeção de subestações em lidar de forma adaptativa e reativa com a complexidade ambiental e a natureza dinâmica dos eventos de interesse}. Falta uma abordagem que integre eficientemente a fusão de múltiplos sensores heterogêneos para criar uma percepção situacional rica e utilize essa percepção para guiar ativamente a navegação, permitindo que o robô seja "atraído" por anomalias relevantes e "repelido" por perigos ou eventos normais, superando os desafios impostos pela \ac{EMI}, obstáculos, iluminação variável e a natureza mista do ambiente. Desenvolver um \textit{framework} que habilite essa capacidade de navegação reativa inteligente é crucial para destravar todo o potencial da inspeção robótica autônoma em subestações elétricas.

\section{Justificativa e Relevância}
A presente pesquisa se justifica pela necessidade premente de aprimorar os métodos de inspeção em subestações elétricas, superando as limitações das abordagens manuais e das atuais soluções robóticas baseadas em rotas. A relevância deste trabalho reside tanto em suas potenciais contribuições práticas para o setor elétrico quanto em seus avanços científicos no campo da robótica autônoma.
Do ponto de vista prático, o desenvolvimento de um \textit{framework} para navegação robótica reativa inteligente tem o potencial de impactar significativamente a operação e manutenção de subestações. Ao permitir inspeções mais frequentes, consistentes e detalhadas, a solução proposta pode levar a uma detecção mais precoce de falhas e anomalias, contribuindo para a redução de custos com manutenção corretiva e minimizando o risco de interrupções no fornecimento de energia \cite{Wang2024Advancements, AtlantisPress2023Analysis}. A automação de tarefas de inspeção, especialmente em ambientes perigosos ou de difícil acesso, aumenta a segurança dos trabalhadores, reduzindo sua exposição a riscos elétricos e ambientais \cite{Anybotics2023Advantages}. Além disso, a capacidade de reação inteligente a eventos inesperados pode otimizar o tempo de inspeção, focando os recursos do robô nas áreas que realmente demandam atenção, aumentando a eficiência geral do processo \cite{EnergyRobotics2024Future}.
Cientificamente, esta tese busca preencher uma lacuna importante no estado da arte da navegação autônoma em ambientes complexos e não estruturados. A proposta de um \textit{framework} que utiliza fusão sensorial heterogênea para gerar sinais de atração/repulsão representa uma abordagem inovadora para a navegação reativa, afastando-se dos paradigmas tradicionais baseados em planejamento de caminho \cite{Hou2023ChallengesSolutions}. A aplicação de aprendizado por reforço para a tomada de decisão adaptativa sobre como reagir a estímulos ambientais complexos (distinguindo eventos relevantes de ruídos ou ocorrências normais) constitui uma contribuição relevante para a inteligência e autonomia robótica \cite{Tai2024AutonomousRL, Kim2023MultimodalFusion}. A investigação sobre a integração eficaz de múltiplos sensores (incluindo \ac{LiDAR}, visão computacional, visão térmica, percepção térmica, áudio, ultrassom) e a sua fusão para criar uma percepção situacional robusta em ambientes com forte \ac{EMI} e obstáculos metálicos também representa um avanço significativo \cite{Hussain2020ReviewChallenges, Kim2024AutonomousNavSensorFusion}. Finalmente, ao abordar os desafios específicos das subestações elétricas, este trabalho contribui com conhecimento aplicável a outros domínios que exigem navegação robótica robusta em condições adversas.
Portanto, a pesquisa proposta é relevante por oferecer uma solução tecnológica com potencial de impacto direto na indústria de energia elétrica, ao mesmo tempo que avança o conhecimento científico em áreas chave da robótica autônoma, como percepção, aprendizado e navegação inteligente em ambientes complexos.

\section{Objetivos}
Com base no problema de pesquisa e na justificativa apresentados, esta tese tem como objetivo principal:

\subsection{Objetivo Geral}
\begin{itemize}
    \item Desenvolver e avaliar um \textit{framework} computacional para navegação robótica reativa inteligente, aplicável à inspeção autônoma de subestações elétricas, que utilize fusão sensorial heterogênea e aprendizado por reforço para gerar sinais de atração e repulsão, guiando o robô de forma dinâmica em resposta a estímulos ambientais.
    
\end{itemize}
Para alcançar este objetivo geral, os seguintes objetivos específicos serão perseguidos:

\subsection{Objetivos Específicos}
\begin{enumerate}
    \item Propor uma arquitetura de \textit{framework} modular: Definir os componentes essenciais do \textit{framework}, incluindo módulos para aquisição e pré-processamento de dados de múltiplos sensores (visuais, térmicos, \ac{LiDAR}, áudio, etc.), fusão sensorial, detecção de eventos, tomada de decisão baseada em aprendizado por reforço (atração/repulsão) e controle de navegação reativa integrada a \ac{SLAM} e modelos 3D.
    
    \item Investigar e implementar técnicas de fusão sensorial: Explorar e selecionar métodos adequados para fundir dados de sensores heterogêneos, visando a criação de uma representação rica e robusta do estado do ambiente da subestação, capaz de identificar potenciais eventos de interesse ou perigo mesmo sob condições desafiadoras (\ac{EMI}, oclusões, iluminação variável).
    
    \item Desenvolver e treinar um modelo de aprendizado por reforço: Aplicar técnicas de aprendizado por reforço, como \textit{Q-Learning} ou variantes, para treinar um agente capaz de decidir autonomamente se um evento detectado pela fusão sensorial deve gerar um sinal de atração (investigar) ou repulsão (evitar/ignorar), com base em recompensas que considerem a relevância da inspeção e a autopreservação do robô.
    
    \item Integrar o \textit{framework} com sistemas de localização e mapeamento: Combinar os sinais de atração/repulsão gerados com algoritmos de \ac{SLAM} e informações de modelos 3D pré-existentes da subestação para gerar trajetórias de navegação reativas, seguras e eficientes, que permitam ao robô investigar anomalias e evitar obstáculos e perigos.
    
    \item Validar o \textit{framework} proposto através de simulação: Avaliar o desempenho e a eficácia do \textit{framework} desenvolvido em um ambiente de simulação que reproduza as características e desafios de uma subestação elétrica, comparando a abordagem reativa com métodos tradicionais baseados em rotas pré-definidas em termos de eficiência de inspeção, robustez da navegação e capacidade de resposta a eventos \cite{collins2021simulation}.
    
\end{enumerate}
\section{Metodologia Geral}
A abordagem metodológica adotada nesta tese será predominantemente construtiva e experimental, focada no desenvolvimento e avaliação do \textit{framework} proposto para navegação robótica reativa inteligente. A pesquisa será conduzida através das seguintes fases principais:
\begin{enumerate}
    \item Revisão Bibliográfica: Será realizada uma revisão sistemática e crítica da literatura científica e técnica relevante, abrangendo as áreas de robótica móvel autônoma, navegação em ambientes complexos, fusão de sensores heterogêneos (\ac{LiDAR}, visão, térmico, áudio), aprendizado por reforço aplicado à robótica, inspeção de subestações elétricas e os desafios específicos associados (EMI, obstáculos, etc.). Esta fase visa consolidar o estado da arte, refinar a identificação das lacunas existentes e fundamentar as escolhas de projeto para o \textit{framework}.
    \item Projeto e Desenvolvimento do \textit{Framework}: Com base na revisão bibliográfica e nos objetivos específicos, será projetada a arquitetura modular do \textit{framework} computacional. Serão definidos os componentes, suas funcionalidades, as interfaces de comunicação entre eles e os algoritmos específicos a serem empregados em cada módulo (por exemplo, técnicas de fusão, algoritmo de aprendizado por reforço, estratégias de controle de navegação reativa). A implementação do \textit{framework} será realizada como um conjunto de pacotes dentro do ecossistema \ac{ROS} (Robot Operating System) ou ROS2, aproveitando sua arquitetura de comunicação baseada em nós, tópicos e serviços, e facilitando a integração com simuladores e hardware robótico padrão. A escolha entre \ac{ROS} ou ROS2 será definida com base nos requisitos específicos e nas ferramentas disponíveis durante o desenvolvimento.
    \item Implementação e Treinamento dos Componentes Chave: Os algoritmos selecionados para fusão sensorial e o modelo de aprendizado por reforço para decisão de atração/repulsão serão implementados como nós ROS/ROS2 e integrados ao \textit{framework}. O agente de aprendizado por reforço será treinado em ambiente simulado (como Gazebo, integrado ao ROS/ROS2), utilizando cenários que representem os desafios e eventos típicos de uma subestação, com definição cuidadosa dos estados, ações e funções de recompensa para guiar o aprendizado em direção ao comportamento desejado (inspeção eficaz e segura).
    \item Validação Experimental por Simulação: A validação do \textit{framework} proposto será conduzida por meio de experimentos em um ambiente de simulação robusto e realista, integrado ao ROS/ROS2 \cite{NVIDIA2024Simulation}. Serão criados cenários de teste que simulem diferentes condições operacionais e desafios encontrados em subestações (presença de EMI simulada, diferentes tipos de obstáculos, variabilidade de iluminação, ocorrência de eventos anômalos). O desempenho do \textit{framework} de navegação reativa será avaliado quantitativa e qualitativamente, utilizando métricas como taxa de sucesso na detecção e investigação de eventos, eficiência da navegação (tempo, distância), robustez a falhas e segurança (evitação de colisões e perigos). Os resultados serão comparados com abordagens de linha de base, como a navegação tradicional baseada em rotas pré-definidas, para demonstrar a superioridade e os benefícios da abordagem reativa proposta \cite{collins2021simulation, Garbacea2022FrameworkVV}.
\end{enumerate}
    
Esta metodologia permitirá o desenvolvimento sistemático do \textit{framework} e uma avaliação rigorosa de seu desempenho em condições controladas, fornecendo evidências para suportar as contribuições científicas e tecnológicas da tese.

\section{Organização do Documento}
Este documento está estruturado da seguinte forma para apresentar a pesquisa de forma clara e lógica:
\begin{itemize}
    \item Capítulo 1 (Introdução): Apresenta a contextualização do tema de inspeção robótica em subestações, define o problema de pesquisa focado nos desafios da navegação autônoma reativa, justifica a relevância prática e científica do trabalho, estabelece os objetivos geral e específicos, descreve a metodologia geral adotada e, finalmente, detalha a organização deste documento.
    \item Capítulo 2 (Revisão da Literatura): Discute em profundidade o estado da arte nas áreas relevantes para esta tese, incluindo robótica móvel autônoma, técnicas de navegação (SLAM, baseada em modelos), fusão de sensores heterogêneos, aprendizado por reforço aplicado à robótica, e as particularidades da inspeção em subestações elétricas. Identifica as lacunas na literatura que esta pesquisa visa preencher.
    \item Capítulo 3 (Proposta do \textit{Framework}): Detalha a arquitetura do \textit{framework} proposto para navegação robótica reativa inteligente. Descreve cada módulo (aquisição/pré-processamento, fusão/detecção, decisão atração/repulsão, navegação reativa), os algoritmos selecionados e as interconexões entre eles.
    \item Capítulo 4 (Implementação e Desenvolvimento): Apresenta os detalhes da implementação do \textit{framework} e dos seus componentes chave, incluindo as ferramentas de software utilizadas, a configuração do ambiente de simulação e o processo de treinamento do agente de aprendizado por reforço.
    \item Capítulo 5 (Resultados e Discussão): Apresenta os resultados dos experimentos de validação realizados em simulação. Analisa o desempenho do \textit{framework} proposto em comparação com abordagens de linha de base, discute as implicações dos resultados, as limitações do estudo e as contribuições alcançadas em relação aos objetivos definidos.
    \item Capítulo 6 (Conclusões e Trabalhos Futuros): Sumariza as principais conclusões da pesquisa, reitera as contribuições científicas e tecnológicas, e aponta direções para trabalhos futuros, incluindo possíveis extensões do \textit{framework} e validação em robôs físicos.
\end{itemize}
